\chapter{Homologia przecięć, snopy perwersyjne i twierdzenie o rozkładzie}
\label{ch:ih-perverse-xiii}

\section{Dlaczego zwykła homologia nie wystarcza dla osobliwości?}
W rozdziale~\ref{ch:snopy-verdier-v12} otrzymaliśmy dualność Verdiera,
ale dla przestrzeni osobliwej kompleks dualizujący nie musi być
przesunięciem snopu stałego. Przykładowo dla węzła $zw=0$
lokalne grupy $H^q(X,X\setminus\{0\};k)$ są niezerowe w dwóch
stopniach, a w punkcie gładkim krzywej jedynie w stopniu $2$.
Podobna trudność pojawia się w dualności Poincarégo: zwykłe grupy
homologii zwartej przestrzeni osobliwej nie muszą układać się w
niedegenerującą parę stopni dopełniających się do wymiaru.

Ideą Goresky'ego--MacPhersona jest ograniczenie łańcuchów, które
mogą przechodzić przez warstwy osobliwe. Liczbę dopuszczalnych
wymiarów przecięcia reguluje funkcja zwana
\emph{perwersyjnością} (ang. \emph{perversity}).
Współczesna konstrukcja snopowa opisuje tę samą teorię przez
\emph{kompleks przecięcia} $IC_X$, który jest szczególnym
przykładem snopu perwersyjnego. Umożliwia to sformułowanie
ważnego twierdzenia o rozkładzie obrazów właściwych.

W części łańcuchowej pracujemy nad ustalonym ciałem $k$,
zazwyczaj $\Q$. W zastosowaniach algebraicznych
przyjmujemy $k=\Q$ i topologię klasyczną zespolonych
rozmaitości algebraicznych. Tam, gdzie twierdzenie wymaga
zaawansowanej teorii pochodnej lub Hodge'a, wskazujemy je
jawnie jako wynik przyjęty, zamiast zastępować dowód samą
formułą.

\section{Pseudorozmaitości warstwowe i funkcje perwersyjności}
\begin{definicja}[Warstwowa pseudorozmaitość PL]
\label{def:pseudomanifold-xiii}
Niech $n\ge2$. Przez $n$-wymiarową pseudorozmaitość
warstwową typu PL rozumiemy parazwartą przestrzeń $X$
z filtracją domkniętymi podzbiorami
$X=X_n\supseteq X_{n-1}=X_{n-2}\supseteq\cdots\supseteq X_0
\supseteq X_{-1}=\varnothing$, zgodną z lokalnymi
triangulacjami, taką że każda różnica
$S_j=X_j\setminus X_{j-1}$ jest (być może pustą)
$j$-rozmaitością PL, a każdy punkt $x\in S_j$ dla $j<n$ posiada
warstwowe otoczenie homeomorficzne z
$\R^j\times cL_x$; tutaj $L_x$ jest zwartą
$(n-j-1)$-wymiarową przestrzenią warstwową,
a $cL_x=(L_x\times[0,1))/(L_x\times\{0\})$
ma naturalną filtrację stożkową. Jest to stożek otwarty:
nie zawiera warstwy bazowej $L_x\times\{1\}$, a po
usunięciu wierzchołka jest homeomorficzny z
$L_x\times(0,1)$. W punktach warstwy
głównej $S_n$ wymagamy otoczeń homeomorficznych
z $\R^n$; nie stosujemy tam wzoru na link
wymiaru $-1$. Dla $j<n$ wymiar linku może wynosić
co najmniej $1$ wobec braku warstw kowymiaru $1$. Zakładamy gęstość
$X\setminus X_{n-2}$ w $X$ i brak warstw kowymiaru $1$.
Przestrzeń nazywamy \emph{normalną}, jeżeli linki
wszystkich warstw osobliwych są spójne. W rozdziale
rozważamy tylko przestrzenie z lokalnie skończoną
triangulacją zgodną z filtracją.
\end{definicja}
\begin{uwaga}
Dla skończonego kompleksu PL (a także dla stratyfikacji
Whitneya zespolonej rozmaitości algebraicznej po przejściu
do odpowiedniego modelu topologicznego) warunki powyższe
zapewniają lokalne modele stożkowe. W ogólnej topologii
istnienie takich modeli nie jest automatyczne.
Dla $S\subset X$ o wymiarze $n-c$ liczbę $c$ nazywamy
\emph{kowymiarem rzeczywistym} warstwy $S$.
\end{uwaga}

\begin{definicja}[Perwersyjność Goresky'ego--MacPhersona]
\label{def:perversity-xiii}
Niech $X$ będzie $n$-wymiarową warstwową
pseudorozmaitością PL, $n\ge2$. Perwersyjnością
klasyczną na $X$ nazywamy funkcję
$\bar p:\{2,\ldots,n\}\to\Z$ spełniającą
$\bar p(2)=0$ oraz
$\bar p(c+1)\in\{\bar p(c),\bar p(c)+1\}$
dla $2\le c<n$. Perwersyjność górna $\bar t$
oraz dwie środkowe dane są przez
\[
\bar t(c)=c-2,\qquad
\bar m(c)=\left\lfloor\frac{c-2}{2}\right\rfloor,
\qquad
\bar n(c)=\left\lceil\frac{c-2}{2}\right\rceil.
\]
Dwie perwersyjności $\bar p,\bar q$ są
\emph{komplementarne}, gdy $\bar p(c)+\bar q(c)=c-2$
dla wszystkich $2\le c\le n$.
\end{definicja}
\begin{przyklad}[Środkowa perwersyjność]
Dla kowymiarów $c=2,3,4,5$ wartości $\bar m(c)$
są kolejno $0,0,1,1$, a wartości $\bar n(c)$
wynoszą $0,1,1,2$. Dla parzystych kowymiarów
$\bar m(c)=\bar n(c)$; dla nieparzystych mogą się
różnić. Ta drobna pozornie rozbieżność będzie istotna
w przykładzie zawieszenia torusa.
\end{przyklad}

\section{Łańcuchy dopuszczalne i operator brzegu}
\begin{definicja}[Dopuszczalność]
\label{def:allowability-xiii}
Niech $X$ będzie $n$-wymiarową pseudorozmaitością PL
z ustaloną triangulacją zgodną z filtracją, $k$ ciałem,
a $\bar p$ klasyczną perwersyjnością na $X$.
Dla łańcucha PL $\xi$ wymiaru $i$
ze współczynnikami w $k$ oznaczmy przez
$|\xi|$ sumę nośników sympleksów o niezerowych
współczynnikach. Łańcuch $\xi$ jest
\emph{$\bar p$-dopuszczalny}, gdy dla każdego $c\ge2$
\[
\dim\bigl(|\xi|\cap X_{n-c}\bigr)
\le i-c+\bar p(c)
\]
oraz te same nierówności zachodzą dla $\partial\xi$
ze stopniem $i-1$. Przyjmujemy
$\dim\varnothing=-\infty$. Dopuszczamy skończone
subdywizje zgodne z filtracją i przechodzimy do
identyfikacji wynikających z subdywizji.
Dla grup $IH_*(cL;k)$ i innych niezwartych przestrzeni
korzystamy z \emph{łańcuchów skończonych}
(o nośniku zwartym), nie z lokalnie skończonych.
Wzory dla dwóch wariantów nie są zamienne.
Kompleks takich łańcuchów oznaczamy
$I^{\bar p}C_*(X;k)$, a jego homologię przez
$IH_i^{\bar p}(X;k)$.
\end{definicja}
\begin{uwaga}[Dlaczego warunek na brzeg jest osobny?]
Nierówność dla $|\xi|$ nie gwarantuje nierówności
dla $|\partial\xi|$: brzeg może leżeć w warstwie
osobliwej, mimo że sam łańcuch ma wyższy wymiar.
Dlatego w definicji jawnie wymagamy dopuszczalności
obu łańcuchów. Wtedy $\partial^2=0$ jest zwykłą
tożsamością symplicjalną i
$I^{\bar p}C_*(X;k)$ rzeczywiście jest kompleksem.
Również na przestrzeniach niezwartych używamy tutaj
łańcuchów o nośniku zwartym;
wersja lokalnie skończona wymaga osobnego warunku
lokalnej skończoności. Nie utożsamiamy tych wersji.
\end{uwaga}
\begin{przyklad}[Łańcuch przechodzący przez osobliwość]
Jeśli $X$ ma wymiar $4$, a punkt $v$ jest warstwą
kowymiaru $4$, to dla $\bar m(4)=1$ i $i=1$
warunek ma postać $\dim(|\xi|\cap\{v\})\le-2$.
Zatem dopuszczalny łańcuch jednowymiarowy nie
przechodzi przez $v$. Dla $i=3$ prawa strona wynosi
$0$, więc spotkanie $v$ przez łańcuch jest dozwolone,
lecz trzeba dodatkowo sprawdzić dopuszczalność jego brzegu.
\end{przyklad}

\section{Wzór na stożek: pierwszy istotny rachunek}
\begin{twierdzenie}[Wzór stożkowy]
\label{thm:cone-ih-xiii}
Niech $L$ będzie zwartą $(n-1)$-wymiarową
warstwową pseudorozmaitością PL, $n\ge2$, a $cL$
jej otwartym stożkiem, którego wierzchołek jest
warstwą kowymiaru $n$. Niech $\bar p$ będzie klasyczną
perwersyjnością, a $k$ ustalonym ciałem. Wówczas
\[
IH_i^{\bar p}(cL;k)\cong
\begin{cases}
IH_i^{\bar p}(L;k),&0\le i<n-1-\bar p(n),\\
0,&i\ge n-1-\bar p(n),
\end{cases}
\]
przy czym po prawej dla $L$ stosujemy perwersyjność
ograniczoną do kowymiarów występujących w $L$.
\end{twierdzenie}
\begin{proof}[Dowód po przyjęciu lematu o dopuszczalnym stożkowaniu PL]
Po usunięciu wierzchołka $v$ mamy
$cL\setminus\{v\}\cong L\times(0,1)$.
Przesunięcie wzdłuż czynnika $(0,1)$ stanowi
warstwową homotopię do poziomu $L\times\{1/2\}$.
Standardowy operator pryzmatu dla tej homotopii
zachowuje nierówności dopuszczalności, gdyż nie
zmienia kowymiaru spotykanej warstwy; w konsekwencji
$IH_i(cL\setminus\{v\};k)\cong IH_i(L;k)$.

Oznaczmy $r=n-1-\bar p(n)$.
Jeżeli $i<r$, dopuszczalny $i$-łańcuch nie może
spotkać $v$, bo nierówność w wierzchołku miałaby
postać $0\le i-n+\bar p(n)<-1$.
Również dopuszczalne homologiczne relacje między
cyklami tego stopnia można zastąpić relacjami poza
wierzchołkiem: ich wymiar $i+1\le r$ nadal wymusza
uniknięcie $v$. Zatem w tym zakresie homologie są
te same co dla $L\times(0,1)$.

Dla $i\ge r$ weźmy dopuszczalny cykl $z$.
Wybierzmy wierzchołek $v$ i utwórzmy geometryczny
stożek $v*z$, po wcześniejszej subdywizji, jeśli
trzeba. Jego brzeg spełnia
$\partial(v*z)=z-v*(\partial z)=z$.
W wierzchołku warunek dopuszczalności dla wymiaru
$i+1$ to $0\le i+1-n+\bar p(n)$, co zachodzi
właśnie dla $i\ge r$. Dla pozostałych warstw
stożkowanie podwyższa jednocześnie wymiar łańcucha
i wymiar przecięcia z daną warstwą co najwyżej
o jeden. Warunki dla $z$ gwarantują więc warunki
dla $v*z$. Jeśli $z$ spotyka $v$, stosujemy następujący
lemat PL: po wspólnej subdywizji łańcuchów w
stożku, triangulacji zgodnej z filtracją oraz
operacji ``stożka łańcuchowego'' istnieje operator
$C_v$, określony na dopuszczalnych cyklach w
stopniu $i\ge r$, taki że
$\partial C_v(z)=z$, a każdy sympleks $C_v(z)$
spełnia warunki dopuszczalności. \emph{Szkic weryfikacji:}
na sympleksach nierozpiętych na $v$ jest to zwykły
stożek; na sympleksach zawierających $v$ stosuje się
normalizację kompleksu symplicjalnego i usuwa
sympleksy zdegenerowane. Wspólna subdywizja
zapewnia, że operatory brzegu oraz te redukcje
są ze sobą zgodne. Warunek w wierzchołku to
$0\le i+1-n+\bar p(n)$, a dla innych warstw
oba wymiary rosną co najwyżej o jeden.
Zatem w wysokich stopniach dopuszczalne cykle
są brzegami.

\emph{Granica dowodu:} zgodność normalizacji z
filtracją oraz pełna konstrukcja operatora $C_v$
dla arbitralnej dopuszczalnej triangulacji są
treścią standardowego lematu o stożkowaniu PL,
który tutaj przyjmujemy. Wzór na próg i
kontrola nierówności są uzasadnione jawnie,
lecz nie traktujemy powyższego szkicu jako
niezależnego dowodu całej teorii łańcuchów
przecięć.
\end{proof}
\begin{przyklad}[Stożek na torusie]
\label{ex:torus-cone-xiii}
Niech $L=T^2$ i $X=cT^2$, więc $\dim_\R X=3$.
Dla $\bar m(3)=0$ wartość graniczna wynosi
$r=3-1-0=2$. Współczynniki $k$ przyjmujemy w ciele.
Z wzoru stożkowego otrzymujemy
$IH_0^{\bar m}(X;k)\cong k$,
$IH_1^{\bar m}(X;k)\cong k^2$ i
$IH_i^{\bar m}(X;k)=0$ dla $i\ge2$.
Dla górnej perwersyjności $\bar n(3)=1$ próg to $r=1$:
pozostaje tylko $IH_0^{\bar n}(X;k)\cong k$.
To pierwszy przykład, w którym dwie środkowe
perwersyjności wykrywają różne dane.
\end{przyklad}

\section{Dualność przecięcia i przykład zawieszenia torusa}
\begin{twierdzenie}[Dualność Goresky'ego--MacPhersona]
\label{thm:ih-duality-xiii}
Niech $X$ będzie zwartą, zorientowaną, normalną
$n$-wymiarową pseudorozmaitością PL bez brzegu,
a $k$ ciałem. Dla komplementarnych perwersyjności
klasycznych $\bar p+\bar q=\bar t$ istnieje
naturalne, niedegenerujące sparowanie
$IH_i^{\bar p}(X;k)\otimes_k
 IH_{n-i}^{\bar q}(X;k)\to k$.
Jeżeli wszystkie występujące warstwy osobliwe mają parzysty
kowymiar rzeczywisty, to dla środkowej
perwersyjności mamy $\bar m=\bar n$ na warstwach,
a więc otrzymujemy samodualność środkowej teorii.
\end{twierdzenie}
\begin{proof}[Zakres przedstawionego argumentu]
Na najwyższej warstwie parowanie powstaje z przecięcia
dopuszczalnych łańcuchów w położeniu ogólnym.
Warunki $\bar p(c)+\bar q(c)=c-2$ gwarantują,
że przecięcia i ich brzegi nie przechodzą przez
warstwy osobliwe w wymiarach, w których mogłyby
zniszczyć zgodność z operatorem brzegu.
Niedegeneracja jest istotą głębokiego twierdzenia
i wymaga indukcji po warstwach, wzoru stożkowego,
lokalnych klas orientacji i odpowiedniej wersji
dualności Verdiera dla kompleksu przecięcia.
Nie uzyskujemy jej z samego liczenia wymiarów:
przyjmujemy ją tutaj jako twierdzenie Goresky'ego--MacPhersona.
\end{proof}

\begin{przyklad}[Dlaczego dolna perwersyjność nie zawsze jest samodualna?]
\label{ex:suspension-torus-xiii}
Niech $X=\Sigma T^2$ będzie zawieszeniem torusa.
Jest to zwarta, zorientowana pseudorozmaitość wymiaru $3$
z dwoma punktami osobliwymi, których linkiem jest $T^2$.
Pokryjmy $X$ dwoma otwartymi stożkami $U,V$; ich
przecięcie $W=U\cap V$ jest warstwowo homotopijnie
równoważne $T^2$. Zatem
$IH_i^{\bar p}(W;k)\cong H_i(T^2;k)$, czyli
w stopniach $0,1,2$ otrzymujemy odpowiednio $k,k^2,k$.

Dla dolnej perwersyjności $\bar m(3)=0$ wzór
stożkowy daje
$IH_0^{\bar m}(U;k)=k$,
$IH_1^{\bar m}(U;k)=k^2$ oraz
$IH_i^{\bar m}(U;k)=0$ dla $i\ge2$;
to samo dotyczy $V$. W długim ciągu
Mayera--Vietorisa mamy, po opuszczeniu
symbolu $\bar m$:
\[
0\longrightarrow IH_3(X)\longrightarrow
H_2(T^2;k)\longrightarrow0,
\]
więc $IH_3(X)\cong k$. Kolejny istotny
fragment ma postać
\[
0\longrightarrow IH_2(X)\longrightarrow
k^2\xrightarrow{\,u\mapsto(u,-u)\,}k^2\oplus k^2
\longrightarrow IH_1(X)\longrightarrow
k\xrightarrow{\,t\mapsto(t,-t)\,}k\oplus k.
\]
Oba zapisane odwzorowania są różnowartościowe.
Pierwsza ma kokernele izomorficzne z $k^2$,
zatem $IH_2(X)=0$ i $IH_1(X)\cong k^2$.
Na końcu tego ciągu otrzymujemy $IH_0(X)\cong k$.

Dla górnej perwersyjności $\bar n(3)=1$ wzór
stożkowy daje dla $U,V$ jedynie grupę
$IH_0^{\bar n}=k$, a wszystkie wyższe znikają.
Wtedy ciąg Mayera--Vietorisa redukuje się
do izomorfizmów
$IH_3^{\bar n}(X)\cong H_2(T^2;k)=k$ oraz
$IH_2^{\bar n}(X)\cong H_1(T^2;k)=k^2$;
w stopniach $1,0$ otrzymujemy dokładny ciąg
\[
0\longrightarrow IH_1^{\bar n}(X)
\longrightarrow k\xrightarrow{\,t\mapsto(t,-t)\,}
k\oplus k\longrightarrow IH_0^{\bar n}(X)
\longrightarrow0.
\]
Stąd $IH_1^{\bar n}(X)=0$ i
$IH_0^{\bar n}(X)\cong k$.
Podsumowując,
\[
IH_i^{\bar m}(X;k)=(k,k^2,0,k),\qquad
IH_i^{\bar n}(X;k)=(k,0,k^2,k)
\quad (i=0,1,2,3).
\]
Niedegenerujące parowanie z
Twierdzenia~\ref{thm:ih-duality-xiii} łączy
$IH_1^{\bar m}$ z $IH_2^{\bar n}$,
a nie z $IH_2^{\bar m}=0$.
Przyczyną jest nieparzysty kowymiar osobliwości;
warunek Witta nie jest tutaj spełniony.
\emph{Zależność dowodowa:} użyliśmy standardowego
ciągu Mayera--Vietorisa dla dopuszczalnych kompleksów
łańcuchów; jego dowód wymaga przejścia do
kompleksu snopowego lub wersji podziału jedności
na łańcuchach PL i jest przyjętym twierdzeniem
pomocniczym tego rozdziału.
\end{przyklad}

\section{Kompleksy konstruktybilne, podpory i podpory nadzwyczajne}
Przenosimy się teraz do zespolonej geometrii
algebraicznej. Niech $X$ będzie \emph{zespoloną odmianą
algebraiczną} skończonego typu (nie zakładamy jej
gładkości) z ustaloną skończoną
stratyfikacją, której warstwy są gładkie i spójne.
Przez $D_c^b(X;k)$ oznaczamy ograniczoną kategorię
pochodną kompleksów snopów $k$-przestrzeni wektorowych,
których snopy kohomologii są lokalnymi systemami
skończonego rzędu na każdej warstwie.
Dla punktu $x\in X$ niech $i_x:\{x\}\hookrightarrow X$.

\begin{definicja}[Warunki podpory i kopodpory]
\label{def:perverse-support-xiii}
Niech $X$ będzie zespoloną odmianą algebraiczną
ze skończoną stratyfikacją na gładkie warstwy,
$k$ ciałem, a $i_x:\{x\}\hookrightarrow X$ włożeniem
punktu $x\in X$. Kompleks $K\in D_c^b(X;k)$ spełnia
\emph{warunek podpory}, jeżeli dla każdego punktu $x$
leżącego w warstwie $S$ wymiaru zespolonego $s$
$H^j(i_x^*K)=0$ dla $j>-s$.
Spełnia \emph{warunek kopodpory}, jeżeli
$H^j(i_x^!K)=0$ dla $j<s$.
Kompleks spełniający oba warunki nazywamy
\emph{snopem perwersyjnym} (w środkowej konwencji
kohomologicznej). Kategorię takich obiektów oznaczamy
$\operatorname{Perv}(X;k)$.
\end{definicja}
\begin{uwaga}[Dlaczego przy kopodporze występuje $i_x^!$?]
Kiełek $i_x^*K$ sprawdza, co widać w sąsiedztwie
punktu $x$, a $i_x^!K$ rejestruje klasy lokalne
\emph{z nośnikiem} w $x$. Występują tu dokładnie
obie operacje omawiane przy dualności Verdiera.
W istocie kopodpora $K$ jest podporą
$\mathbb D_XK$ po odpowiedniej dualizacji.
\end{uwaga}

\begin{stwierdzenie}[Podstawowy przykład: rozmaitość gładka]
\label{prop:smooth-perverse-xiii}
Niech $M$ będzie gładką zespoloną rozmaitością
wymiaru zespolonego $d$, $k$ ciałem, a $L$ lokalnym
systemem skończeniewymiarowych $k$-przestrzeni wektorowych
na $M$. Wtedy $L[d]$ jest snopem perwersyjnym.
\end{stwierdzenie}
\begin{proof}
Dla punktu $x\in M$ zwykły kiełek $i_x^*L[d]$
ma jedyną niezerową kohomologię w stopniu $-d$.
Z drugiej strony lokalna kohomologia pary
$(B^{2d},B^{2d}\setminus\{0\})$ ma jedyną niezerową
grupę w stopniu $2d$. Dlatego $i_x^!L[d]$
ma jedyną niezerową kohomologię w stopniu
$2d-d=d$. Obie nierówności z
Definicji~\ref{def:perverse-support-xiii}
są spełnione. Ponieważ konstrukcja zachowuje
lokalne trywializacje $L$, wynik nie zależy od
wybranego punktu.
\end{proof}

\begin{twierdzenie}[Serce struktury $t$]
\label{thm:heart-perverse-xiii}
Dla zespolonej rozmaitości algebraicznej $X$
i ciała $k$ warunki podpory i kopodpory wyznaczają
ograniczoną strukturę $t$ na $D_c^b(X;k)$.
Jej sercem jest abelowa kategoria
$\operatorname{Perv}(X;k)$; istnieją funktory
kohomologii perwersyjnej $\,{}^p\!H^j$.
Dualność Verdiera zachowuje to serce i zamienia
warunek podpory z warunkiem kopodpory.
\end{twierdzenie}
\begin{proof}[Status dowodu]
Pełny dowód wymaga konstrukcji perwersyjnych
truncacji, sklejenia struktur $t$ wzdłuż podziału
na warstwy oraz sprawdzenia dokładności funktorów
odnośnych do włożenia otwartego i domkniętego.
Nie ma go w poprzednim rozdziale, dlatego
przyjmujemy tutaj twierdzenie o strukturze $t$
jako jawny wynik podstawowy. Nierówności w
Definicji~\ref{def:perverse-support-xiii}
są pełną definicją testowaną na punktach,
a nie skrótowym określeniem „złego stopnia”.
\end{proof}

\section{Przedłużenie pośrednie i kompleks przecięcia}
\begin{definicja}[Przedłużenie pośrednie]
\label{def:middle-extension-xiii}
Niech $X$ będzie zespoloną odmianą algebraiczną,
$k$ ciałem, $j:U\hookrightarrow X$ otwartym włożeniem,
a $P\in\operatorname{Perv}(U;k)$.
Przedłużeniem pośrednim $j_{!*}P$ nazywamy obraz
morfizmu w kategorii abelowej $\operatorname{Perv}(X;k)$:
\[
{}^p\!H^0(Rj_!P)\longrightarrow {}^p\!H^0(Rj_*P).
\]
Jest to jedyny perwersyjny kompleks przedłużający $P$,
który nie ma nietrywialnego podobiektu ani ilorazu
podpartego w $X\setminus U$.
\end{definicja}
\begin{definicja}[Kompleks przecięcia]
\label{def:ic-xiii}
Niech $X$ będzie zredukowaną, czystowymiarową zespoloną
odmianą algebraiczną wymiaru zespolonego $d$,
$k$ ciałem, a $j:U=X_{\mathrm{reg}}\hookrightarrow X$ włożeniem
gładkiej części. Jej \emph{kompleksem przecięcia}
ze współczynnikami w $k$ nazywamy
$IC_X(k)=j_{!*}(k_U[d])$.
Dla dowolnego $i\in\Z$ definiujemy
\[
IH^i(X;k):=\mathbb H^{i-d}(X,IC_X(k)),
\]
gdzie $\mathbb H$ oznacza hiperkohomologię kompleksu snopów.
Dla gładkiego $X$ otrzymujemy $IC_X(k)=k_X[d]$
i $IH^i(X;k)=H^i(X;k)$.
\end{definicja}
\begin{uwaga}[Dwie konwencje stopni]
Często autorzy oznaczają przez $IC_X$ kompleks
\emph{nieprzesunięty}, którego hiperkohomologia
jest $IH^i$. U nas $IC_X$ jest \emph{perwersyjny}
i zawiera przesunięcie $[d]$; stąd konieczne jest
$\mathbb H^{i-d}$, a nie $\mathbb H^i$.
Nie należy mylić tego z przesunięciem rzeczywistego
wymiaru $2d$ występującym w dualności Verdiera.
\end{uwaga}

\begin{twierdzenie}[Porównanie wersji łańcuchowej i snopowej]
\label{thm:ic-comparison-xiii}
Niech $X$ będzie zwartą zespoloną odmianą
algebraiczną czystego wymiaru $d$,
stratyfikowaną zgodnie z lokalnym modelem stożkowym,
a współczynniki należą do ciała $k$.
Wówczas hiperkohomologia perwersyjnego kompleksu
$IC_X(k)$ oblicza środkową kohomologię przecięć
Goresky'ego--MacPhersona, z uwzględnieniem
konwencji przesunięcia w Definicji~\ref{def:ic-xiii}.
Ponadto $\mathbb D_XIC_X(k)\simeq IC_X(k)$,
przy standardowej orientacji zespolonej części gładkiej
i perwersyjnej konwencji przesunięcia o $[d]$.
\end{twierdzenie}
\begin{proof}[Zakres]
Podstawą porównania jest lokalny wzór stożkowy
i aksjomaty charakteryzujące kompleks przecięcia:
warunek normalizacji na części gładkiej,
odpowiednie ograniczenia podpór oraz reguła
przedłużenia przez kolejne warstwy osobliwe.
Dowód przeprowadza się indukcyjnie po stratach;
nie wynika on automatycznie z definicji $j_{!*}$.
Przyjmujemy go jako twierdzenie porównawcze.
Samodualność korzysta z komplementarności
środkowych perwersyjności, które na przestrzeniach
zespolonych zgadzają się na kowymiarach warstw
\emph{zespolonych} (czyli parzystych kowymiarach
rzeczywistych). Nie twierdzimy tego dla dowolnej
stratyfikacji rzeczywistej; w szczególności
przykład $\Sigma T^2$ nie jest modelem
algebraicznym o tej własności.
\end{proof}

\section{Obliczenie: krzywa węzłowa i jej normalizacja}
\label{sec:nodal-perverse-xiii}
Rozpatrzmy najpierw lokalnie
$X=\{zw=0\}\cap B_\varepsilon\subset\C^2$.
Normalizacja $\nu:\widetilde X\to X$ rozdziela
obie gałęzie; $\widetilde X$ jest sumą rozłączną
dwóch małych dysków. Jest to odwzorowanie skończone,
a zatem właściwe, o zerowym wymiarze włókien.

\begin{stwierdzenie}[Kompleks przecięcia węzła]
\label{prop:node-ic-xiii}
Niech $B_\varepsilon\subset\C^2$ będzie kulą o środku $0$,
$X=\{(z,w)\in B_\varepsilon:zw=0\}$, a
$\nu:\widetilde X\to X$ normalizacją.
Dla dowolnego ciała $k$ zachodzi naturalny quasi-izomorfizm
$IC_X(k)\simeq\nu_*k_{\widetilde X}[1]$.
W szczególności w węźle $x=0$ mamy
$H^{-1}(i_x^*IC_X(k))\cong k^2$,
a wszystkie pozostałe grupy kohomologii kiełka
są zerowe.
\end{stwierdzenie}
\begin{proof}
Ponieważ $\nu$ jest skończona, jej włókna są
skończone i $R^q\nu_*k_{\widetilde X}=0$ dla $q>0$;
w naszym lokalnym modelu wynika to również
bezpośrednio z faktu, że przeciwobrazy dostatecznie
małych otoczeń są sumą dwóch dysków.
Kompleks $K=\nu_*k_{\widetilde X}[1]$ ogranicza się
do $k_U[1]$ na części gładkiej $U$.
W punkcie $x$ jego zwykły kiełek ma kohomologię
$k^2$ wyłącznie w stopniu $-1$. Natomiast
$H^q(i_x^!K)$ obliczamy przez lokalne pary
dysków i ich nakłutych wersji: niezerowa grupa
$k^2$ pojawia się wyłącznie dla $q=1$.
Takie dane spełniają warunki podpory i kopodpory;
w szczególności $H^0(i_x^*K)=H^0(i_x^!K)=0$.
Przy użyciu adjunkcji dla włożenia punktu
ostatnie dwa zaniki oznaczają, że $K$ nie ma
podobiektów ani ilorazów perwersyjnych
podpartych wyłącznie w $x$. Z jednoznaczności
przedłużenia pośredniego otrzymujemy tezę.
\end{proof}

\begin{przyklad}[Krzywa rzutowa: kohomologia zwykła i przecięć]
Niech $C$ będzie nieredukowalną zespoloną
krzywą rzutową z jednym zwyczajnym węzłem,
której normalizacja $\nu:\mathbb P^1\to C$
jest izomorfizmem poza dwiema różnymi preobrazami
węzła (np. odpowiednia wymierna krzywa sześcienna).
Topologicznie $C$ powstaje z $S^2$ przez
utożsamienie dwóch punktów, więc
$H^0(C;k)\cong H^2(C;k)\cong k$,
$H^1(C;k)\cong k$.
Z drugiej strony
$IC_C(k)=\nu_*k_{\mathbb P^1}[1]$ i właściwość $\nu$
dają $IH^i(C;k)\cong H^i(\mathbb P^1;k)$.
A zatem $IH^1(C;k)=0$ przy zachowaniu
$IH^0(C;k)=IH^2(C;k)=k$.
Nowa teoria nie „ignoruje” osobliwości: rozdziela
lokalne gałęzie i odzyskuje prawidłową dualność.
\end{przyklad}

\begin{przyklad}[Kohomologia krzywej węzłowej z dokładnego ciągu snopów]
\label{prz:node-normalization-exact-sequence}
W sytuacji poprzedniego przykładu oznaczmy
przez $p\in C$ węzeł, a przez
$p_+,p_-\in\mathbb P^1$ dwa punkty jego przeciwobrazu.
Mamy następujący \emph{krótki ciąg dokładny snopów}
wektorowych na $C$:
\[
 0\longrightarrow k_C
 \xrightarrow{\ \eta\ }\nu_*k_{\mathbb P^1}
 \xrightarrow{\ \Delta_p\ }k_{\{p\}}
 \longrightarrow0.
\]
Pierwsza strzałka przypisuje lokalnie stałej funkcji
na $C$ jej złożenie z normalizacją.
Druga przypisuje parze lokalnych wartości na obu
gałęziach ich różnicę:
$\Delta_p(a_+,a_-)=a_+-a_-$.
Sprawdźmy dokładność na kiełkach, korzystając
z Lematu~\ref{lem:stalk-exact-v12}.
Poza węzłem $\nu$ jest izomorfizmem, a ciąg na
kiełkach redukuje się do
$0\to k\xrightarrow{\mathrm{id}}k\to0$.
W węźle otrzymujemy natomiast
\[
 0\longrightarrow k
 \xrightarrow{\ a\mapsto(a,a)\ }k^2
 \xrightarrow{\ (a_+,a_-)\mapsto a_+-a_-\ }k
 \longrightarrow0,
\]
co jest dokładne także dla ciała charakterystyki $2$.

Ponieważ $\nu$ ma skończone włókna i lokalne
przeciwobrazy małych otoczeń są rozłącznymi dyskami,
$R^q\nu_*k_{\mathbb P^1}=0$ dla $q>0$.
Wobec tego
$H^q(C,\nu_*k_{\mathbb P^1})
\cong H^q(\mathbb P^1;k)$.
Z długiego ciągu kohomologii otrzymujemy fragment
\[
 0\to H^0(C;k)\to k
 \xrightarrow{\ 0\ }k
 \to H^1(C;k)\to0
\]
oraz izomorfizm
$H^2(C;k)\cong H^2(\mathbb P^1;k)=k$.
Środkowa strzałka jest zerowa: globalna funkcja
lokalnie stała na spójnej $\mathbb P^1$ ma tę samą
wartość w $p_+$ i $p_-$.
Dlatego
\[
 H^0(C;k)=k,\qquad
 H^1(C;k)=k,\qquad
 H^2(C;k)=k.
\]
To \emph{niezależny rachunek snopowy} potwierdzający
wcześniejsze rozumowanie o sklejeniu dwóch punktów.

Z kolei $IC_C(k)=\nu_*k_{\mathbb P^1}[1]$
na mocy Stwierdzenia~\ref{prop:node-ic-xiii}.
Przy naszej konwencji otrzymujemy
\[
\mathbb H^{i-1}(C,IC_C(k))=H^i(\mathbb P^1;k).
\]
Wobec tego
$IH^0(C;k)=IH^2(C;k)=k$ i
$IH^1(C;k)=0$.
Rozbieżność w stopniu pierwszym wynika dokładnie
z ilorazu $k^2/\operatorname{diag}(k)$
w węźle, a nie z niejawnej zmiany współczynników.
\end{przyklad}

\begin{figure}[!htbp]
\centering
\begin{tikzpicture}[>=Stealth,font=\small]
 \node[draw,rounded corners=4pt,fill=manifoldblue!8,
       minimum width=2.65cm,minimum height=1.1cm,
       align=center] (a) at (0,0)
       {$k$\\
        funkcja na $C$};
 \node[draw,rounded corners=4pt,fill=manifoldgreen!8,
       minimum width=2.65cm,minimum height=1.1cm,
       align=center] (b) at (4.1,0)
       {$k^2$\\
        dwie gałęzie};
 \node[draw,rounded corners=4pt,fill=manifoldgold!12,
       minimum width=2.65cm,minimum height=1.1cm,
       align=center] (c) at (8.2,0)
       {$k$\\
        różnica wartości};
 \draw[->,thick] (a)--(b)
       node[midway,above=6pt] {$\eta_p$};
 \draw[->,thick] (b)--(c)
       node[midway,above=6pt] {$\Delta_p$};
 \node[align=center] at (4.1,-1.45)
       {$H^1(C;k)\cong k$,\qquad
        $IH^1(C;k)=0$};
\end{tikzpicture}
\caption{Kiełkowy krótki ciąg dokładny normalizacji
krzywej węzłowej. Iloraz dwóch wartości gałęzi
przez ich diagonalną identyfikację wyjaśnia
dodatkową klasę zwykłej kohomologii w stopniu $1$.}
\label{fig:node-normalization-exact-sequence}
\end{figure}

\begin{center}
\begin{tikzpicture}[scale=0.9,font=\small,>=Stealth]
\draw[manifoldblue,thick] (-3.8,0) ellipse (0.62 and 0.74);
\draw[manifoldblue,thick] (-2.3,0) ellipse (0.62 and 0.74);
\fill[manifoldred] (-3.18,0) circle (2.3pt);
\fill[manifoldred] (-2.92,0) circle (2.3pt);
\node[above left] at (-3.18,0.08) {$p_+$};
\node[above right] at (-2.92,0.08) {$p_-$};
\node[above] at (-3.1,1.0) {dwie gałęzie};
\draw[->,thick] (-1.35,0)--(0.7,0)
 node[midway,above] {$\nu$};
\draw[manifoldblue,thick] (1.65,0) ellipse (1.02 and 0.8);
\draw[manifoldblue,thick] (2.2,0.1) ellipse (1.02 and 0.8);
\fill[manifoldred] (1.92,0.71) circle (2.3pt);
\node[above] at (2.3,1.0) {węzeł};
\node[below] at (-3, -1.1) {$\widetilde X$ (schematycznie)};
\node[below] at (2, -1.1) {$X$ (schematycznie)};
\end{tikzpicture}
\par\smallskip\small Normalizacja oddziela lokalne gałęzie; rysunek
pokazuje ideę, a nie dokładny obraz algebraiczny krzywej.
\end{center}

\begin{przyklad}[Snop stały na krzywej węzłowej nie jest przedłużeniem pośrednim]
\label{prz:node-constant-vs-IC-24}
Niech $X=\{zw=0\}\cap B_\varepsilon$ z węzłem $p=0$.
Zwróćmy uwagę, że perwersyjny kompleks $k_X[1]$
\emph{nie jest} izomorficzny z $IC_X(k)$, mimo że
oba ograniczają się do $k[1]$ na części gładkiej.
Rzeczywiście, rachunek z
Przykładu~\ref{ex:node-duality-v12} daje
\[
 H^j(i_p^!k_X[1])\cong H^{j+1}(X,X\setminus\{p\};k)
 =\begin{cases}
 k,&j=0,\\ k^2,&j=1,\\0,&j\notin\{0,1\}.
 \end{cases}
\]
Z kolei zwykły kiełek $i_p^*k_X[1]$ ma jedyną
niezerową kohomologię $k$ w stopniu $-1$.
Wraz z identycznym rachunkiem na warstwie gładkiej
pokazuje to, że $k_X[1]$ spełnia obie nierówności
perwersyjności (w punkcie $p$ wymiar warstwy wynosi $0$).
Jednak $H^0(i_p^!k_X[1])=k$ daje przez adjunkcję
niezerowy morfizm snopów perwersyjnych
$k_{\{p\}}\to k_X[1]$. Ponieważ snop punktowy
$k_{\{p\}}$ jest prosty, jest to monomorfizm:
$k_X[1]$ ma niezerowy podobiekt podparty w $p$.
Tymczasem dla $IC_X(k)=\nu_*k_{\widetilde X}[1]$
Stwierdzenie~\ref{prop:node-ic-xiii} daje
$H^0(i_p^!IC_X(k))=0$, więc taki podobiekt
nie istnieje. Widać wprost, dlaczego warunek
\emph{przedłużenia pośredniego} jest silniejszy
od samej perwersyjności, a zwykły snop stały
nie zastępuje kompleksu przecięcia na osobliwości.
\end{przyklad}

\section{Twierdzenie o rozkładzie i małe odwzorowania}
\begin{definicja}[Odwzorowanie małe i półmałe]
Niech $f:Y\to X$ będzie właściwym morfizmem
zespolonych odmian algebraicznych czystego wymiaru $d$,
generycznie skończonym i dominującym. Dla $r\ge1$
niech $X_r=\{x\in X:\dim_\C f^{-1}(x)\ge r\}$.
Odwzorowanie jest \emph{półmałe}, jeśli
$\operatorname{codim}_\C X_r\ge2r$ dla każdego $r\ge1$,
a \emph{małe}, jeśli wszędzie zachodzi
nierówność ostra. Przyjmujemy
$\operatorname{codim}\varnothing=+\infty$.
\end{definicja}
\begin{uwaga}
Przy powyższej definicji już założyliśmy, że $f$ jest
dominujący i generycznie skończony pomiędzy
czystowymiarowymi przestrzeniami tego samego wymiaru;
dla ogólnych morfizmów półmałość formułuje się przez
stratyfikację i nierówność
$2\dim f^{-1}(x)\le\dim Y-\dim S$
dla $x$ w warstwie $S$.
\end{uwaga}

\begin{twierdzenie}[Twierdzenie o rozkładzie --- wersja projektowa]
\label{thm:decomposition-xiii}
Niech $f:Y\to X$ będzie projektowym morfizmem
zespolonych odmian algebraicznych, a $Y$
niech będzie gładka i czystowymiarowa o wymiarze
zespolonym $d$. W kategorii $D_c^b(X;\Q)$ istnieje
(niekanoniczny na ogół) izomorfizm
\[
Rf_*\Q_Y[d]\simeq
\bigoplus_{j\in\Z}{}^p\!H^j(Rf_*\Q_Y[d])[-j].
\]
Każdy snop perwersyjny po prawej jest półprosty,
a więc rozkłada się jako skończona suma kompleksów
$IC_{\overline S}(L)$, przedłużonych przez zero
do $X$, gdzie $S$ przebiega odpowiednie gładkie,
lokalnie domknięte warstwy, a $L$ są
półprostymi lokalnymi systemami $\Q$-przestrzeni
wektorowych na $S$. W konsekwencji
$Rf_*\Q_Y[d]$ rozkłada się na przesunięte
kompleksy przecięcia.
\end{twierdzenie}
\begin{proof}[Jawnie używany wynik głęboki]
Nie dowodzimy tu twierdzenia o rozkładzie BBDG.
Jego klasyczne dowody wymagają czystości i wag,
polaryzowanych struktur Hodge'a albo innych
silnych twierdzeń geometrycznych.
Z konstrukcji perwersyjnej struktury $t$ wynika
istnienie snopów ${}^p\!H^j$, ale \emph{nie}
wynika z niej rozszczepienie kompleksu na sumę
swoich kohomologii perwersyjnych.
Właśnie to rozszczepienie i półprostota są
nietrywialną treścią przyjętego twierdzenia.
\end{proof}
\begin{wniosek}[Rozdzielczości małe]
\label{cor:small-map-ic-xiii}
Niech $f:Y\to X$ będzie właściwą \emph{małą}
rozdzielczością osobliwości, tj. biracjonalnym
morfizmem z gładkiej czystowymiarowej
$Y$ na nieredukowalną $X$, który jest izomorfizmem
nad gęstą częścią gładką, przy $d=\dim_\C X=\dim_\C Y$.
Zakładamy, że dla $X_r=\{x\in X:\dim_\C f^{-1}(x)\ge r\}$
zachodzi $\operatorname{codim}_\C X_r>2r$ dla każdego
$r\ge1$, z konwencją $\operatorname{codim}\varnothing=+\infty$.
Wówczas $Rf_*\Q_Y[d]\simeq IC_X(\Q)$.
\end{wniosek}
\begin{proof}[Wyjaśnienie zastosowanego kryterium]
Małość eliminuje nietrywialne składniki
perwersyjne podparte w mniejszych stratach;
na części otwartej $Rf_*\Q_Y[d]$ ogranicza się
do $\Q_U[d]$. Przyjmujemy standardowe twierdzenie
o małych mapach (wynikające z aksjomatów
kompleksu przecięcia), które daje jego
identyfikację z przedłużeniem pośrednim.
Małość nie oznacza jedynie, że włókna są spójne:
potrzebne są jawne nierówności wymiarów.
\end{proof}

\begin{przyklad}[Normalizacja krzywej węzłowej jest odwzorowaniem małym]
\label{prz:small-node-normalization}
Powróćmy do nieredukowalnej rzutowej krzywej węzłowej
$C$ o normalizacji
$\nu:\mathbb P^1\to C$ z
Przykładu~\ref{prz:node-normalization-exact-sequence}.
Morfizm $\nu$ jest skończony, stąd właściwy,
a każde jego włókno jest zbiorem skończonym:
poza węzłem składa się z jednego punktu,
a nad węzłem z dwóch punktów.
Dla każdego $r\ge1$ zbiór
$C_r=\{x\in C:\dim_\C\nu^{-1}(x)\ge r\}$
jest więc pusty. Nierówność
$\operatorname{codim}_\C C_r>2r$
zachodzi z przyjętą konwencją
$\operatorname{codim}\varnothing=+\infty$.
Ponadto normalizacja jest izomorfizmem
nad gęstym otwartym podzbiorem
$C\setminus\{p\}$ i ma gładką dziedzinę
$\mathbb P^1$.
Spełnia zatem \emph{wszystkie} hipotezy
Wniosku~\ref{cor:small-map-ic-xiii} przy $d=1$
dla $k=\Q$. Dla dowolnego ciała $k$ to samo
wynika z bezpośredniego rachunku lokalnego
ze Stwierdzenia~\ref{prop:node-ic-xiii}:
\[
 R\nu_*k_{\mathbb P^1}[1]\simeq IC_C(k).
\]
W tym szczególnym przykładzie izomorfizm
sprawdziliśmy już niezależnie na kiełkach
i kopiełkach (Stwierdzenie~\ref{prop:node-ic-xiii}).
Z kolei rozdmuchanie początku $\C^2$ opisane niżej
\emph{nie} jest małe: jego jednowymiarowe włókno
ma kowymiar dokładnie $2$, a wymagana byłaby
nierówność ostra $2>2$.
To zestawienie odróżnia małość od półmałości
na dwóch jawnych morfizmach.
\end{przyklad}

\section{Przykład geometryczny: rozdmuchanie płaszczyzny}
\label{sec:blowup-decomp-xiii}
Niech $f:\widetilde{\C^2}\to\C^2$ będzie
rozdmuchaniem początku. Jest to projektowe,
a więc właściwe odwzorowanie; poza $0$ jest
izomorfizmem, natomiast $f^{-1}(0)\cong\mathbb P^1$.
W szczególności $f$ jest półmałe, lecz nie małe:
$\operatorname{codim}_\C\{0\}=2=2\dim_\C f^{-1}(0)$.

\begin{stwierdzenie}[Rozkład dla rozdmuchania]
\label{prop:blowup-decomp-xiii}
Niech $f:\widetilde{\C^2}\to\C^2$ będzie rozdmuchaniem
początku. W kategorii $D_c^b(\C^2;\Q)$ mamy
\[
Rf_*\Q_{\widetilde{\C^2}}[2]
\simeq\Q_{\C^2}[2]\oplus\Q_{\{0\}}.
\]
Dla rozdmuchania punktu na $\mathbb P^2$ otrzymujemy
$\dim_\Q H^0=\dim_\Q H^4=1$ oraz
$\dim_\Q H^2=2$ dla rozdmuchanej powierzchni.
\end{stwierdzenie}
\begin{proof}
Twierdzenie o zmianie bazy dla właściwego odwzorowania
utożsamia kiełek $R^qf_*\Q$ w punkcie $x$
z $H^q(f^{-1}(x);\Q)$. Dla $x\ne0$ włókno jest
punktem, a dla $x=0$ jest $\mathbb P^1$.
Wynika stąd $R^0f_*\Q\cong\Q_{\C^2}$,
$R^2f_*\Q\cong\Q_{\{0\}}$ i zanik pozostałych.
Nie wolno z tego \emph{samego} wnioskować, że
$Rf_*\Q$ jest sumą swoich snopów kohomologii:
morfizmy łączące w kategorii pochodnej mogą
być nietrywialne.

Ponieważ $f$ jest półmałe, stosujemy standardowe
kryterium półmałości: $P=Rf_*\Q_{\widetilde{\C^2}}[2]$
jest snopem perwersyjnym. Twierdzenie
\ref{thm:decomposition-xiii} zapewnia jego półprostotę.
Na $U=\C^2\setminus\{0\}$ ograniczenie $P$ jest
równe $\Q_U[2]$, toteż jedynym składnikiem
prostym o pełnym nośniku jest
$IC_{\C^2}(\Q)=\Q_{\C^2}[2]$,
z krotnością jeden. Pozostałe składniki proste
mogą być podparte tylko w $\{0\}$, więc
$P\simeq\Q_{\C^2}[2]\oplus\Q_{\{0\}}^{\oplus m}$.
Krotność wyznaczamy z kiełka w początku:
$H^0(i_0^*P)\cong H^2(\mathbb P^1;\Q)\cong\Q$,
podczas gdy $H^0(i_0^*\Q_{\C^2}[2])=0$.
Zatem $m=1$, co uzasadnia rozszczepienie z tezy.
Po zastosowaniu $R\Gamma$ do analogicznego wzoru
dla rozdmuchania $\mathbb P^2$ dodatkowa klasa
w stopniu $2$ zwiększa drugą liczbę Bettiego o jeden,
przy zachowaniu zerowego i czwartego stopnia.
\end{proof}
\begin{uwaga}[Skąd potrzeba półprostoty?]
Dane o wszystkich $R^qf_*\Q$ nie określają
jeszcze klasy izomorfizmu $Rf_*\Q$ w kategorii
pochodnej. Twierdzenie o rozkładzie kontroluje
również klasy rozszerzeń między tymi składnikami.
Właśnie dlatego jest tak silnym narzędziem
w geometrii algebraicznej.
\end{uwaga}

\section{Względne twarde twierdzenie Lefschetza}
\begin{twierdzenie}[Względne twarde Lefschetz]
\label{thm:relative-hard-lefschetz-xiii}
Niech $f:Y\to X$ będzie projektowym morfizmem
zespolonych odmian algebraicznych, $Y$
gładką czystowymiarową o wymiarze $d$,
a $\eta\in H^2(Y;\Q)$ klasą pierwszego Cherna
wiązki liniowej względnie ample względem $f$.
Wtedy dla każdego $r\ge0$ mnożenie przez
$\eta^r$ wyznacza izomorfizm
\[
\eta^r:{}^p\!H^{-r}(Rf_*\Q_Y[d])
\longrightarrow{}^p\!H^{r}(Rf_*\Q_Y[d]).
\]
\end{twierdzenie}
\begin{proof}[Status]
Jest to niezależne głębokie twierdzenie, a nie
formalny wniosek z definicji snopu perwersyjnego.
Pełny dowód wymaga aparatu Hodge'a lub teorii wag.
Warto jednak zauważyć analogię ze zwykłym twardym
twierdzeniem Lefschetza dla gładkich odmian
projektowych: tu stopnie kohomologii zastąpione
są stopniami perwersyjnymi.
\end{proof}
\begin{przyklad}[Projekcja z włóknem $\mathbb P^1$]
Niech $f:\mathbb P^1\times\mathbb P^1\to\mathbb P^1$
będzie pierwszą projekcją. Wtedy
$Rf_*\Q_{\mathbb P^1\times\mathbb P^1}[2]
\simeq\Q_{\mathbb P^1}[2]\oplus\Q_{\mathbb P^1}$.
Ponieważ $P=\Q_{\mathbb P^1}[1]$ jest perwersyjny,
dwie składowe są $P[1]$ i $P[-1]$, zatem
${}^p\!H^{-1}\cong P\cong{}^p\!H^{1}$.
Klasa $\eta$ pochodząca z drugiego czynnika
przenosi generator $H^0(\mathbb P^1;\Q)$ na
generator $H^2(\mathbb P^1;\Q)$ w każdym włóknie;
stąd w tym przypadku izomorfizm względnego
Lefschetza można sprawdzić bezpośrednio.
\end{przyklad}

\section{Granice aktualnego rozdziału i pomost do geometrii algebraicznej}
Uzupełniliśmy schemat
\[
\text{łańcuchy dopuszczalne}
\longrightarrow IH^*(X)
\longleftrightarrow IC_X
\longrightarrow\operatorname{Perv}(X)
\longrightarrow Rf_*IC_Y.
\]
Łańcuchowa konstrukcja i rachunki stożkowe
mają szczegółowe uzasadnienia. Pojęcie snopu
perwersyjnego jest zdefiniowane przez
sprawdzalne warunki na kiełkach i kopiełkach;
przykłady rozmaitości gładkiej oraz węzła
zostały obliczone. Twierdzenie o rozkładzie,
względne twarde Lefschetz i pełna równoważność
konstrukcji łańcuchowej ze snopową są
\emph{jawnie przyjętymi} głębokimi twierdzeniami.

Przy późniejszym scalaniu z geometrią algebraiczną
warto dołączyć rozdziały o strukturach Hodge'a,
monodromii i cyklach zanikających, dualności
Grothendiecka oraz o \emph{$t$-strukturach} w kategoriach
pochodnych. Nie utożsamiamy tu kohomologii
singularnej i przecięć ze snopami koherentnymi
ani z kohomologią etalną bez dodatkowych twierdzeń.



\section{Zakres uzasadnień i przyjętych twierdzeń}
\label{sec:xiii-status-po-audycie}
\begin{uwaga}[Zakres samodzielności dowodów]
Rachunek wzoru stożkowego wymaga standardowego lematu
PL o stożkowaniu dopuszczalnych łańcuchów; jawny
rachunek nierówności nie zastępuje konstrukcji
tych operatorów. Formuła Mayera--Vietorisa dla
homologii przecięć wymaga odpowiedniej wersji
kompleksów łańcuchów dopuszczalnych i ich
niezmienniczości przy przejściu do pokryć otwartych.
Z kolei porównanie $IC_X$ z teorią łańcuchową,
istnienie perwersyjnej struktury $t$, dualność
Goresky'ego--MacPhersona, twierdzenie o rozkładzie
i względne twarde Lefschetz są \emph{jawnie
przyjętymi} głębokimi wynikami. W dalszym
rozdziale rozwijamy aparat kategorii pochodnych
i teorii Hodge'a potrzebny do precyzyjnego
sformułowania tych zależności; nie ogłaszamy
przez to automatycznie pełnych dowodów tych twierdzeń.
\end{uwaga}
